%
% menu_lifecycle_delivery.tex
%
% A refinement of docs/menu_lifecycle.tex along the ONE axis that model's
% delivery abstraction collapses: the fate of an enqueued-but-undrained menu
% event across a session departure. It exists to adjudicate a round-6 review
% challenge (PR #495, bead lux-m3xr) to the M1 "no false-positive delivery"
% invariant.
%
% THE CHALLENGE. The base model's MenuDeliver is one atomic step guarded on
% (registered AND hasInbox): it fuses MenuEventRouter.deliver's live-set read
% and inbox.offer's get-then-put, and it carries no token for the message once
% enqueued. So it cannot represent, and therefore cannot rule on, the ordering
% the reviewer named:
%
%   t0  deliver() reads live_sessions(): the target c is live. (registry lock,
%       released)
%   t1  a departure removes c from the client registry.
%   t2  deliver() calls offer(). offer wins _inboxes_lock, the inbox is STILL
%       present, it puts the event and returns True -> deliver reports
%       "delivered".
%   t3  the departure cascade's inbox drop_session runs, discarding that inbox
%       WITH the just-enqueued event still in it.
%
% This is offer-THEN-drop-discards, distinct from the base model's D1
% (drop-BETWEEN-get-and-put, an orphan put). The D1 fix -- offer's get and put
% under one _inboxes_lock hold -- does NOT close t0..t3: the discard at t3 is a
% separate lock acquisition that legitimately follows a fully-completed offer.
% No queue can forbid its recipient from departing after a delivery.
%
% THE QUESTION this model answers. Is t0..t3 a false delivery? It is a false
% delivery only if the discarded event had a LIVE recipient that could have
% drained it. So this model gives the enqueued event a token (pending), a drain
% (recv), and -- crucially -- a TWO-PHASE departure whose registry-discard and
% inbox-teardown are separate steps, exactly the "separate critical sections"
% the reviewer says the base model fused. It then checks:
%
%   M1b -- no live-recipient loss:  NOT (lostLive = fset).
%     No undrained menu event is discarded from a connection that is still in
%     the registry. Equivalently: a discard can only ever befall a connection
%     that has already departed, which has no recv() to miss.
%
% WHY t0..t3 IS BENIGN, and why the model proves it. The code's departure
% (hub_display.py _depart / _depart_lapsed) removes the connection from the
% registry (self._clients.discard) and only THEN runs the cascade tail
% (self._cascade.run -> inbox.drop_session), under one StoreLock hold; the
% cascade's own docstring states "connection_id has already been irreversibly
% removed from the registry by the time this runs". So at t3 the recipient is
% already gone from live_sessions(): the discarded event had no live drain to
% lose. DepartRegistry-then-DepartInbox encodes that order; M1b is NOT found.
%
% The reviewer's stated consequence -- "will not re-push the menu" -- does not
% occur either: the SAME departure fires MenuDepartureSink (operations/menus.py),
% which prunes the owner's bar and calls mark_menus() to re-push, on every
% trigger. The click-dispatch path's "delivered" report is irrelevant to the
% re-push; the departure re-pushes regardless. This model carries menuOwner and
% prunes it in DepartRegistry (the same step that deregisters) to show the
% Hub-side half of that fact (menuOwner subset of registered, MO); the display
% re-push is the F3 mechanism, outside this Hub-state model's scope, as
% rendering is throughout menu_lifecycle.tex.
%
% The fidelity control menu_lifecycle_delivery_reorder_buggy.tex runs the two
% departure legs in the WRONG order (inbox teardown before registry discard) and
% makes M1b's negation FOUND -- so the discard-before-cascade order the code
% guarantees is shown to be load-bearing, not assumed.
%
% Type-check:  fuzz docs/menu_lifecycle_delivery.tex
% Model-check:
%   probcli docs/menu_lifecycle_delivery.tex -model_check -p DEFAULT_SETSIZE 3
% Safety goal (NOT found in this, the fixed, spec):
%   probcli docs/menu_lifecycle_delivery.tex -model_check -nodead \
%     -goal "lostLive = fset" -p DEFAULT_SETSIZE 3
%

\documentclass[a4paper,10pt,fleqn]{article}
\usepackage[margin=1in]{geometry}
\usepackage{fuzz}
\usepackage[colorlinks=true,linkcolor=blue,citecolor=blue,urlcolor=blue]{hyperref}

\begin{document}

\title{lux Agent-Menu Delivery: the Undrained-Event Refinement}
\author{Proving no menu click reported delivered is lost to a still-live
  recipient --- the offer-then-drop-discards question the base model abstracts}
\date{September 2026}
\maketitle

% ============================================================
\section{Introduction}
% ============================================================

This specification refines \texttt{menu\_lifecycle.tex} along one axis: the
life of a menu event \emph{after} \texttt{offer} has enqueued it. The base
model's \texttt{MenuDeliver} reports \texttt{delivered} and stops; it carries
no token for the enqueued message, models departure as a single atomic step,
and so cannot state what becomes of a click that was put onto a live inbox
which is then torn down by a departure. A round-6 review of PR~\#495 asked
precisely that: if \texttt{MenuEventRouter.deliver}'s live-set read passes,
then a departure removes the connection, then \texttt{offer} wins its lock and
puts (returning \texttt{True}), and only then the cascade's
\texttt{drop\_session} discards the inbox with the event still in it --- is the
\texttt{delivered} report a false positive?

The base model's D1 addresses a different window: a \texttt{drop\_session}
\emph{between} \texttt{offer}'s \texttt{get} and its \texttt{put} (an orphan
put), closed by holding \texttt{\_inboxes\_lock} across both. The window here is
\emph{after} a fully-completed \texttt{offer}. No lock can prevent a recipient
from departing once it has been delivered to; the honest question is whether
the loss matters, and the honest answer is a property about \emph{who} loses.

\begin{description}
  \item[M1b --- no live-recipient loss.] No undrained menu event is discarded
    from a connection still present in the registry:
    \[ \lnot\, (lostLive = fset). \]
    A click reported \texttt{delivered} may still be discarded --- but only
    when its recipient has already departed, and a departed recipient has no
    \texttt{recv()} left to drain it. The loss is then not a lost service but
    the correct disposal of a click for a session that is gone.
\end{description}

As in \texttt{menu\_lifecycle.tex}, M1b is kept out of the state schema and
checked as the reachability of its negation: \texttt{lostLive = fset} is
\emph{not found} against this fixed spec, and \emph{FOUND} against the
reorder control.

% ============================================================
\section{Basic Types}
% ============================================================

\begin{zed}
  [CONN]
\end{zed}

A connection identity, the \texttt{ConnectionId} keying the client registry,
the inbox table, and the menu registry. Connection ids are unique and never
reused; \texttt{Connect} below encodes that a departing id cannot reconnect
mid-departure. Two \texttt{CONN} suffice to exhibit one owner departing while a
delivery is mid-flight beside a second session; the carrier is bounded by
ProB's \texttt{DEFAULT\_SETSIZE}.

% ============================================================
\section{Free Types}
% ============================================================

\begin{zed}
  DSTATE ::= dnone | ddelivered | dgone
\end{zed}

The outcome \texttt{MenuEventRouter.deliver} last reported to its caller.

\begin{zed}
  FSTATE ::= fclear | fset
\end{zed}

A two-value flag. \texttt{lostLive} carries it: whether any undrained event has
been discarded from a still-registered connection.

\begin{zed}
  TPHASE ::= tidle | tinbox
\end{zed}

Where a departure is in its two-leg cascade: \texttt{tidle} between departures,
\texttt{tinbox} after the registry-and-ownership leg has run and before the
inbox-teardown leg does. This is the interleaving point the base model's atomic
departure hides.

% ============================================================
\section{State}
% ============================================================

\begin{schema}{MenuFlow}
  registered : \power CONN \\
  hasInbox : \power CONN \\
  menuOwner : \power CONN \\
  pending : \power CONN \\
  reported : DSTATE \\
  lostLive : FSTATE \\
  departPhase : TPHASE \\
  departFor : \power CONN
\where
  \# departFor \leq 1
\end{schema}

\texttt{registered} is the client registry's live set --- what
\texttt{live\_sessions()} returns. \texttt{hasInbox} is the set with a live
inbox queue. \texttt{menuOwner} is the set with a stored bar.
\texttt{pending} is the set whose inbox holds an enqueued-but-undrained menu
event; this is the token the base model lacks. \texttt{reported} is the last
\texttt{deliver} outcome. \texttt{lostLive} records a live-recipient loss.
\texttt{departPhase} and \texttt{departFor} (the at-most-one connection whose
departure cascade is in flight) split the departure into its two code legs.

No containment is assumed: \texttt{menuOwner} $\subseteq$ \texttt{registered}
is MO, proved in the base model, and \texttt{hasInbox} is unconstrained so a
delivered-to inbox can be torn down from under a completed \texttt{offer}.

% ============================================================
\section{Initialization}
% ============================================================

\begin{schema}{Init}
  MenuFlow'
\where
  registered' = \emptyset \\
  hasInbox' = \emptyset \\
  menuOwner' = \emptyset \\
  pending' = \emptyset \\
  reported' = dnone \\
  lostLive' = fclear \\
  departPhase' = tidle \\
  departFor' = \emptyset
\end{schema}

% ============================================================
\section{Operations}
% ============================================================

\subsection{Connect --- arrival and reconnect}

A session joins the registry. A departing connection id cannot reconnect while
its cascade is in flight (\texttt{c?} $\notin$ \texttt{departFor}), because
connection ids are unique --- a reconnect is a fresh id, never the one being
retired. This forbids the one spurious interleaving the two-phase departure
would otherwise admit: a same-id ``reconnect'' re-registering between the
registry leg and the inbox leg, which would make the inbox teardown appear to
strike a live connection.

\begin{schema}{Connect}
  \Delta MenuFlow \\
  c? : CONN
\where
  c? \notin departFor \\
  registered' = registered \cup \{ c? \} \\
  hasInbox' = hasInbox \\
  menuOwner' = menuOwner \\
  pending' = pending \\
  reported' = reported \\
  lostLive' = lostLive \\
  departPhase' = departPhase \\
  departFor' = departFor
\end{schema}

\subsection{Arm --- bind the inbox writer}

Models \texttt{inbox.ensure\_writer}: a registered session gains a live inbox.

\begin{schema}{Arm}
  \Delta MenuFlow \\
  c? : CONN
\where
  c? \in registered \\
  hasInbox' = hasInbox \cup \{ c? \} \\
  registered' = registered \\
  menuOwner' = menuOwner \\
  pending' = pending \\
  reported' = reported \\
  lostLive' = lostLive \\
  departPhase' = departPhase \\
  departFor' = departFor
\end{schema}

\subsection{Own --- store a menu bar}

Models \texttt{MenuOperations.set\_menu} for an admitted caller, atomic as the
base model's \texttt{SetMenu} (identity gating and the D3 store window are the
base model's concern). A registered session gains an inbox and a stored bar.

\begin{schema}{Own}
  \Delta MenuFlow \\
  c? : CONN
\where
  c? \in registered \\
  hasInbox' = hasInbox \cup \{ c? \} \\
  menuOwner' = menuOwner \cup \{ c? \} \\
  registered' = registered \\
  pending' = pending \\
  reported' = reported \\
  lostLive' = lostLive \\
  departPhase' = departPhase \\
  departFor' = departFor
\end{schema}

\subsection{MenuDeliver --- the offer leg, guarded on the inbox alone}
\label{sec:menudeliver}

Models \texttt{offer} succeeding: the inbox is present, the message is put, and
the caller is told \texttt{delivered}. The guard is \texttt{c?} $\in$
\texttt{hasInbox} \emph{only}, not \texttt{registered}: this is the faithful
part the base model fuses away. \texttt{MenuEventRouter.deliver} checks the
live set first and then calls \texttt{offer}, which re-checks only the inbox;
by the time \texttt{offer} runs the live-set read may be stale, so a delivery
into the inbox of a connection already removed from the registry (the reviewer's
$t_2$) is exactly a reachable \texttt{MenuDeliver} here. The enqueued event
becomes \texttt{pending}. Any \texttt{hasInbox} connection was registered when
its inbox was armed, so no \texttt{MenuDeliver} fires that the router could not
have reached; guarding on the inbox alone over-approximates delivery, which
only strengthens an M1b that holds.

\begin{schema}{MenuDeliver}
  \Delta MenuFlow \\
  c? : CONN
\where
  c? \in hasInbox \\
  pending' = pending \cup \{ c? \} \\
  reported' = ddelivered \\
  registered' = registered \\
  hasInbox' = hasInbox \\
  menuOwner' = menuOwner \\
  lostLive' = lostLive \\
  departPhase' = departPhase \\
  departFor' = departFor
\end{schema}

\subsection{MenuProviderGone --- the honest non-delivery}

Models \texttt{offer} finding no inbox, and equally the router's earlier
live-set read failing: either way nothing is put and the caller is told
\texttt{provider\_gone}.

\begin{schema}{MenuProviderGone}
  \Delta MenuFlow \\
  c? : CONN
\where
  c? \notin hasInbox \\
  reported' = dgone \\
  registered' = registered \\
  hasInbox' = hasInbox \\
  menuOwner' = menuOwner \\
  pending' = pending \\
  lostLive' = lostLive \\
  departPhase' = departPhase \\
  departFor' = departFor
\end{schema}

\subsection{Drain --- recv consumes the event}

Models the owning session draining its inbox with \texttt{recv()}: a live inbox
with a pending event has it consumed. This is the honest fate of a delivered
click and keeps M1b non-vacuous --- delivered events are not \emph{always}
discarded.

\begin{schema}{Drain}
  \Delta MenuFlow \\
  c? : CONN
\where
  c? \in hasInbox \\
  c? \in pending \\
  pending' = pending \setminus \{ c? \} \\
  registered' = registered \\
  hasInbox' = hasInbox \\
  menuOwner' = menuOwner \\
  reported' = reported \\
  lostLive' = lostLive \\
  departPhase' = departPhase \\
  departFor' = departFor
\end{schema}

\subsection{DepartRegistry --- the registry-and-ownership leg (discard first)}
\label{sec:departregistry}

Models the first half of \texttt{HubDisplay.\_depart}: under \texttt{StoreLock},
\texttt{self.\_clients.discard(c?)} removes the connection from the registry and
ownership is released, \emph{before} the cascade tail runs. The connection
leaves \texttt{registered} and \texttt{menuOwner} in the same step --- the
Hub-side of the ghost-menu re-push (MenuDepartureSink's prune) --- and the
cascade enters its inbox-teardown leg (\texttt{tinbox}), capturing \texttt{c?}.
The inbox is deliberately \emph{not} yet torn down: that is the window in which
a stale-read \texttt{MenuDeliver} can still put onto \texttt{c?}'s live inbox.

\begin{schema}{DepartRegistry}
  \Delta MenuFlow \\
  c? : CONN
\where
  c? \in registered \\
  departPhase = tidle \\
  registered' = registered \setminus \{ c? \} \\
  menuOwner' = menuOwner \setminus \{ c? \} \\
  departPhase' = tinbox \\
  departFor' = \{ c? \} \\
  hasInbox' = hasInbox \\
  pending' = pending \\
  reported' = reported \\
  lostLive' = lostLive
\end{schema}

\subsection{DepartInbox --- the cascade inbox-teardown leg (after the discard)}
\label{sec:departinbox}

Models the cascade tail: \texttt{inbox.drop\_session(c?)} tears down the inbox
and discards any undrained event. It runs only after \texttt{DepartRegistry}
(\texttt{departPhase} = \texttt{tinbox}, \texttt{c?} $\in$ \texttt{departFor}),
so by the code's own guarantee \texttt{c?} has already left \texttt{registered}.
The \texttt{lostLive} guard fires only if a discarded pending event belonged to
a \emph{still-registered} connection --- which cannot happen here, since
\texttt{DepartRegistry} removed \texttt{c?} and \texttt{Connect} cannot re-add
it while its cascade is in flight. So \texttt{lostLive} stays \texttt{fclear}:
the discard is the correct disposal of a click for a departed session.

\begin{schema}{DepartInbox}
  \Delta MenuFlow \\
  c? : CONN
\where
  departPhase = tinbox \\
  c? \in departFor \\
  hasInbox' = hasInbox \setminus \{ c? \} \\
  pending' = pending \setminus \{ c? \} \\
  (c? \in pending \land c? \in registered \implies lostLive' = fset) \\
  (\lnot (c? \in pending \land c? \in registered) \implies lostLive' = lostLive) \\
  departPhase' = tidle \\
  departFor' = \emptyset \\
  registered' = registered \\
  menuOwner' = menuOwner \\
  reported' = reported
\end{schema}

% ============================================================
\section{Verification}
\label{sec:verify}
% ============================================================

Type-checking is a \texttt{fuzz} pass:
\begin{verbatim}
  fuzz docs/menu_lifecycle_delivery.tex
\end{verbatim}

M1b is checked as the reachability of its negation. \emph{Not found} means no
reachable state discards an undrained event from a still-registered connection.

\begin{verbatim}
  # M1b (goal NOT found in this, the fixed, spec)
  probcli docs/menu_lifecycle_delivery.tex -model_check -nodead \
    -goal "lostLive = fset" \
    -p DEFAULT_SETSIZE 3
\end{verbatim}

The reviewer's own interleaving is confirmed \emph{reachable} and benign: a
click reported \texttt{delivered} whose recipient then fully departs, with
nothing lost to a live session.

\begin{verbatim}
  # offer-then-drop-discards, benign (goal FOUND):
  #   Connect(c); Arm(c); MenuDeliver(c); DepartRegistry(c); DepartInbox(c)
  #   leaves reported=ddelivered, registered empty, lostLive=fclear.
  probcli docs/menu_lifecycle_delivery.tex -model_check -nodead \
    -goal "reported = ddelivered & registered = {} & lostLive = fclear" \
    -p DEFAULT_SETSIZE 3

  # the happy path -- delivered AND drained to a live recipient (goal FOUND)
  probcli docs/menu_lifecycle_delivery.tex -model_check -nodead \
    -goal "reported = ddelivered & pending = {} & not(registered = {})" \
    -p DEFAULT_SETSIZE 3
\end{verbatim}

Deadlock-freedom and the structural invariant over the reachable state space:

\begin{verbatim}
  probcli docs/menu_lifecycle_delivery.tex -model_check -p DEFAULT_SETSIZE 3
\end{verbatim}

The reorder control makes M1b's negation \emph{FOUND}:

\begin{verbatim}
  # cascade legs in the wrong order -> live-recipient loss (goal FOUND)
  probcli docs/menu_lifecycle_delivery_reorder_buggy.tex -model_check -nodead \
    -goal "lostLive = fset" \
    -p DEFAULT_SETSIZE 3
\end{verbatim}

Re-run \texttt{fuzz} and these goal checks whenever
\texttt{menu\_event.py}'s \texttt{deliver}, \texttt{inbox.py}'s
\texttt{offer}/\texttt{drop\_session}, \texttt{operations/menus.py}'s
\texttt{MenuDepartureSink}, or \texttt{hub\_display.py}'s
\texttt{\_depart}/\texttt{\_depart\_lapsed} (the discard-before-cascade order)
change.

\end{document}
